%==============================================================
%  Khaled Lakhdher — Curriculum Vitae (Français)
%  Compilation :  pdflatex khaled_cv_fr.tex
%  Fonctionne sur Overleaf, MiKTeX et TeX Live (paquets standard).
%==============================================================
\documentclass[10pt,a4paper]{article}

\usepackage[T1]{fontenc}
\usepackage[utf8]{inputenc}
\usepackage[french]{babel}
\usepackage{charter}                 % police serif proche de l'original
\usepackage[margin=1.25cm,top=0.8cm,bottom=0.7cm]{geometry}
\usepackage[dvipsnames]{xcolor}
\usepackage{titlesec}
\usepackage{enumitem}
\usepackage{hyperref}

% ---------- Couleurs ----------
\definecolor{accent}{HTML}{2C3E63}   % bleu marine des titres
\definecolor{muted}{HTML}{52606D}    % gris texte secondaire

\hypersetup{colorlinks=true,urlcolor=accent,linkcolor=accent}

% ---------- Style des sections ----------
\titleformat{\section}
  {\color{accent}\bfseries\large\scshape}
  {}{0em}{}[{\color{muted!40}\titlerule[0.8pt]}]
\titlespacing*{\section}{0pt}{5pt}{2pt}

% ---------- Style des listes ----------
\setlist[itemize]{leftmargin=1.3em,itemsep=0.5pt,topsep=1pt,parsep=0pt,label=\textbullet}

\pagestyle{empty}
\setlength{\parindent}{0pt}
\linespread{0.96}

% ---------- Aides ----------
\newcommand{\entry}[2]{\textbf{#1}\hfill{\small\color{muted}#2}\par}
\newcommand{\org}[1]{{\itshape\color{muted}#1}\par}
\newcommand{\proj}[1]{\textbf{#1}\par}

%==============================================================
\begin{document}

% ---------- En-tête ----------
\begin{center}
  {\LARGE\bfseries\color{accent} Khaled Lakhdher}\\[3pt]
  {\large\color{muted} Ingénieur en Informatique \;|\; IA \& Data Science \;|\; Développeur IA Full-Stack}\\[5pt]
  {\small\color{muted}
    Sousse, Tunisie \quad|\quad
    \href{mailto:ing.khaledlakhdher@gmail.com}{ing.khaledlakhdher@gmail.com} \quad|\quad
    +216 29 843 899 \quad|\quad
    \href{https://www.linkedin.com/in/khaled-lakhdher}{in/khaled-lakhdher}
  }
\end{center}

% ---------- Profil ----------
\section{Profil}
Étudiant passionné en ingénierie du Machine Learning, avec une solide base en deep learning, vision par
ordinateur et NLP. Maîtrise de Python et TensorFlow et expérience pratique du développement et du
déploiement de modèles d'IA, motivé à contribuer à des solutions innovantes en intelligence artificielle.

% ---------- Expérience ----------
\section{Expérience}
\entry{Ingénieur IA --- Projet de Fin d'Études (PFE)}{Févr. 2026 -- Juin 2026 \textperiodcentered\ Tunisie}
\org{CTT -- Carthage Training \& Technologies}
\begin{itemize}
  \item Conception et développement de \textbf{Hotel Mind}, un assistant conversationnel de réservation hôtelière sur une architecture RAG ancrée connectée à une API SOAP en temps réel, avec une couche de mitigation des hallucinations vérifiant chaque affirmation factuelle.
  \item Livraison d'une solution full-stack conteneurisée atteignant un nDCG@5 de 0,82, un MRR@5 de 0,90 et un taux de réussite de 100\,\% sur un jeu de requêtes de référence.
\end{itemize}
\vspace{2pt}
\entry{Stagiaire Développeur en Intelligence Artificielle}{Juil. 2025 -- Sept. 2025 \textperiodcentered\ Hammamet}
\org{CTT -- Carthage Training \& Technologies}
\begin{itemize}
  \item Fine-tuning d'un modèle NLP Hugging Face pour classifier les avis clients par sentiment (positif, négatif, neutre).
  \item Développement d'un backend Django pour gérer les données clients et automatiser l'export des résultats vers Excel pour le reporting.
  \item Réalisation de l'ensemble des scripts de prétraitement, d'entraînement et d'intégration en Python.
\end{itemize}

% ---------- Projets ----------
\section{Projets}

\proj{Nexus AI --- Plateforme de Gestion d'Employés IA}
\begin{itemize}
  \item Plateforme multi-tenant (FastAPI, PostgreSQL/pgvector, Redis, Next.js) pilotant des agents LLM avec
        permissions IAM, mémoire vectorielle, orchestration multi-agents (bus Redis) et supervision temps réel
        par WebSocket, via une passerelle de modèles agnostique (Anthropic, OpenRouter, Gemini).
\end{itemize}

\proj{Assistant Vocal Personnalisé}
\begin{itemize}
  \item Fine-tuning de Whisper (STT) et SpeechT5 (TTS) avec clonage de voix via CSM-1B pour créer un
        assistant vocal personnalisé, réduisant le temps de réponse de 30\,\%.
\end{itemize}

\proj{Colorisation d'Images par IA \& Génération Automatique de Code}
\begin{itemize}
  \item Entraînement d'un GAN Pix2Pix (U-Net / PatchGAN) pour la colorisation automatique d'images et
        développement d'un système multi-agents (Planner / Architect / Coder) générant des projets
        full-stack (Flask, HTML/CSS/JS, API) à partir d'un simple prompt, servis via une interface web interactive.
\end{itemize}

\proj{Portfolio de Développeur assisté par IA}
\begin{itemize}
  \item Développement d'un portfolio responsive avec React, TypeScript et Tailwind, intégrant un assistant
        RAG (LLM via OpenRouter ancré sur mon CV avec une couche anti-hallucination), déployé sur Vercel.
\end{itemize}

% ---------- Certifications ----------
\section{Certifications}
\begin{itemize}
  \item Fundamentals of Deep Learning --- NVIDIA {\color{muted}(22 nov. 2022)}
  \item Building Transformer-Based Natural Language Processing Applications --- NVIDIA {\color{muted}(19 déc. 2025)}
  \item Applications of AI for Predictive Maintenance --- NVIDIA {\color{muted}(27 oct. 2025)}
\end{itemize}

% ---------- Compétences ----------
\section{Compétences}
\textbf{\color{accent}Langages de programmation :} Python \textperiodcentered\ Java \textperiodcentered\ C \textperiodcentered\ SQL \textperiodcentered\ JavaScript \textperiodcentered\ TypeScript \textperiodcentered\ React.js \textperiodcentered\ React Native \textperiodcentered\ HTML/CSS \textperiodcentered\ Tailwind CSS\par
\smallskip
\textbf{\color{accent}Bibliothèques \& Frameworks IA :} Transformers \textperiodcentered\ PyTorch \textperiodcentered\ TensorFlow \textperiodcentered\ Keras \textperiodcentered\ Hugging Face \textperiodcentered\ Hugging Face Hub \textperiodcentered\ scikit-learn \textperiodcentered\ Pandas \textperiodcentered\ NumPy \textperiodcentered\ Matplotlib\par
\smallskip
\textbf{\color{accent}Frameworks \& Outils de Développement :} Gradio \textperiodcentered\ Streamlit \textperiodcentered\ Flask \textperiodcentered\ Django \textperiodcentered\ Zeep \textperiodcentered\ LangSmith \textperiodcentered\ Pydantic \textperiodcentered\ Docker Compose \textperiodcentered\ Nginx / Reverse Proxy \textperiodcentered\ Postman \textperiodcentered\ Git \textperiodcentered\ GitHub\par
\smallskip
\textbf{\color{accent}Concepts \& Technologies :} Intelligence Artificielle (IA) \textperiodcentered\ Machine Learning (ML) \textperiodcentered\ Deep Learning (DL) \textperiodcentered\ Large Language Models (LLMs) \textperiodcentered\ Agents IA \textperiodcentered\ IA Agentique \textperiodcentered\ Tool-Calling \textperiodcentered\ Prompt Engineering \textperiodcentered\ NLP \textperiodcentered\ ASR \textperiodcentered\ GANs \textperiodcentered\ Bases de Données Vectorielles \textperiodcentered\ Recherche Sémantique \textperiodcentered\ Embeddings \textperiodcentered\ Recherche d'Information \textperiodcentered\ Développement d'API \textperiodcentered\ NoSQL \textperiodcentered\ Cloud Computing\par

% ---------- Formation ----------
\section{Formation}
\entry{Ingénieur en Informatique --- Spécialité IA \& Data Science}{2023 -- 2026}
\org{IEPI -- École Polytechnique Internationale \textperiodcentered\ Sousse, Tunisie}
\vspace{2pt}

\entry{Cycle Préparatoire}{2021 -- 2023}
\org{EPI -- École Polytechnique Internationale Privée \textperiodcentered\ Sousse, Tunisie}

% ---------- Langues ----------
\section{Langues}
\textbf{\color{accent}Français :} DELF B2 \qquad
\textbf{\color{accent}Anglais :} TOEIC (Listening \& Reading) --- 825/990

\end{document}
